\documentclass[10pt,a4paper]{amsart}
\usepackage[margin=19mm]{geometry}
\usepackage{amsmath,amssymb,mathtools}
\usepackage[scr=rsfs,cal=euler]{mathalfa}
\usepackage{xfrac}
\usepackage[unicode,hidelinks,bookmarks=false]{hyperref}
\newcommand{\ie}{i.e.\ }
\newcommand{\cf}{cf.\ }
\newcommand{\resp}{respectively\ }
\newcommand{\Subsection}{\subsection}
\providecommand{\nobreakdash}{\nobreak\hbox{-}}
\setlength{\emergencystretch}{2em}
\multlinegap0pt
\usepackage{bm}
\usepackage{bbm}
\usepackage{dsfont}
\usepackage{xspace}
\usepackage{cancel}
\usepackage{mathtools}
\usepackage{amsthm}
\makeatletter
\@ifundefined{theorem}{\newtheorem{theorem}{Theorem}}{}
\@ifundefined{lemma}{\newtheorem{lemma}[theorem]{Lemma}}{}
\makeatother
\title[On isometric asymptotes of operators quasisimilar to isometr]{On isometric asymptotes of operators quasisimilar to isometries}
\author{Maria F. Gamal'}
\date{Source: arXiv:2201.01752v1 (CC BY 4.0)}
\begin{document}
\maketitle

\noindent\emph{Source and licence.} On isometric asymptotes of operators quasisimilar to isometries. Version arXiv:2201.01752v1 (CC BY 4.0), distributed under the Creative Commons Attribution 4.0 International licence, as recorded in the arXiv licence metadata for this exact version and shown on the abstract page https://arxiv.org/abs/2201.01752v1. Full rights notice: licenses/arxiv-2201.01752-CC-BY-4.0.txt.\par\smallskip

\noindent\textit{Editorial scope.} Counted unit: the Lemma whose source label is \texttt{lemssmain} in the article (source0.tex lines 896--907, source label \texttt{lemssmain}), delivered together with the article's own complete original proof of it (source0.tex lines 909--974, one \texttt{proof} environment of 66 lines). No mathematics is written, repaired or re-derived: statement and proof are the article's own text line for line. Required context retained uncounted: y1 (lines 819--822). Every definition, display and label of those lines is retained unchanged. Omitted from this package and not redistributed: 1--818, 823--895, 908--908, 975--1323. Nothing inside a retained excerpt is omitted or altered: every retained body is delivered exactly as the article's lines are written, so no omission receipt of any kind accompanies this package. Citations: the retained excerpts carry no citation command at all, so no bibliography entry of the article is transcribed and no reference is redistributed. Reference adaptation: none was needed and none was made; every reference inside the delivered unit resolves inside it, so no unresolved reference or citation is produced. Printed slips kept literal: no printed word, symbol, hypothesis or number of the counted statement or of its proof was altered. Layout and class: the article's own class is \texttt{amsart} with packages including hyperref, amsmath, amsfonts, amssymb, amsthm, babel, enumerate, graphicx, color, amsrefs; the wrapper is the lane's pinned \texttt{amsart} editorial wrapper at 10pt on A4, which restates the theorem-like environments the retained text needs and adds the article's own macro definitions that the retained text needs (none), with extra wrapper packages (bm, bbm, dsfont, xspace, cancel, mathtools). The delivered numbering is the unit's own. Assets: no retained span contains a figure environment, a graphics-inclusion command, a PDF-page inclusion, a table or a TikZ picture; no image file is copied into this package, the package declares no asset dependency and no image file is redistributed with it. Licence: version arXiv:2201.01752v1 of this article is distributed under the Creative Commons Attribution 4.0 International licence, as recorded in the arXiv licence metadata for this exact version and shown on the abstract page https://arxiv.org/abs/2201.01752v1. A full rights notice is delivered at licenses/arxiv-2201.01752-CC-BY-4.0.txt. Characters: every retained excerpt is pure ASCII, so the delivered engine, XeLaTeX with its default Latin Modern text font, reports no missing character.
\begin{equation}
\label{y1}
k_{{\mathbf u},\zeta}(z)=\frac{1-\overline{{\mathbf u}(\zeta)}{\mathbf u}(z)}{1-\overline\zeta z},\ \ \ z\in\mathbb D,  
\end{equation}

\begin{lemma}\label{lemssmain} Suppose that $\mathbf v$ is an inner function. Set  
\begin{equation}\label{eev} E_{\mathbf v}=\{\zeta\in\mathbb T\ : \ k_{\mathbf v,\zeta}\in H^2\},
\end{equation}
where $k_{\mathbf v,\zeta}$ are defined by \eqref{y1} for $\mathbf v$. 
Let $m(E_{\mathbf v})=1$. Suppose that $\varphi_0\in H^\infty$ is outer, 
 $1/\varphi_0\not\in H^2$,
 and $\varphi_1\in H^2$ is such that $\varphi_0\varphi_1\in H^\infty$ and 
$\varphi_1\mathcal K_{\mathbf v}\subset H^2$.
Then there exists $\varphi\in H^2$ such that 
$$\varphi\varphi_0\in H^\infty, \ \ \ \varphi\mathcal K_{\mathbf v}\subset H^2, \ \ \ 
 \ \text{ and } \ \ \frac{\varphi}{\varphi_1}\not\in H^\infty.$$
\end{lemma}

\begin{proof} Note that $\mathop{\text{\rm ess\,inf}}_{\mathbb T}|\varphi_0||\varphi_1|=0$. 
Indeed, if there exists $\delta>0$ such that $|\varphi_0||\varphi_1|\geq \delta$  $m$-a.e. on $\mathbb T$, 
then $|\varphi_1|\geq \delta/|\varphi_0|$  $m$-a.e. on $\mathbb T$. Since $\varphi_0$ is outer, 
the latest estimate implies that $1/\varphi_0\in H^2$, which contradicts to the assumption on $\varphi_0$. 

Take $\{C_n\}_{n=1}^\infty$ such that $0<C_n<C_{n+1}$ for all $n\geq 1$ and $C_n\to\infty$. 
Set $$E_n=\{\zeta\in E_{\mathbf v}\ :\ C_n<\|k_{\mathbf v,\zeta}\|\leq C_{n+1}\}, \ \ n \geq 1,$$
and $$\delta_n =\mathop{\text{\rm ess\,inf}}_{E_n}|\varphi_0||\varphi_1|,\ \ n \geq 1.$$
(If $m(E_n)=0$, then $\delta_n=+\infty$.)

Consider two cases. \emph{First case}: there exists $n\geq 1$ such that $\delta_n=0$. Then 
there exist sequences $\{\delta_{1k}\}_{k=1}^\infty$ and $\{\tau_{1k}\}_{k=1}^\infty$ such that
 $0<\delta_{1k+1}<\delta_{1k}$, $\delta_{1k}\to 0$, $\tau_{1k}\subset E_n$, $m(\tau_{1k})>0$, 
and $$\delta_{1k+1}<|\varphi_0||\varphi_1|\leq\delta_{1k} \  m\text{-a.e. on }\tau_{1k} \text{ for all }k\geq 1. $$
Set $A_k=1/\delta_{1k}$, $k\geq 1$. Take a sequence $\{\varepsilon_k\}_{k=1}^\infty$ such that 
$\varepsilon_k>0$  for all $k\geq 1$, and 
\begin{equation}\label{aaepsilon} \sum_{k=1}^\infty\varepsilon_kA_k^2<\infty.\end{equation}
There  exists a sequence $\{\tau_k\}_{k=1}^\infty$ such that $\tau_k\subset\tau_{1k}$, $m(\tau_k)>0$, 
and $$\int_{\tau_k}|\varphi_1|^2{\mathrm d}m\leq\varepsilon_k \ \text{ for all }k\geq 1.$$
Then \begin{equation}\label{aavarphi1}  \sum_{k=1}^\infty A_k^2\int_{\tau_k}|\varphi_1|^2{\mathrm d}m<\infty\end{equation}
and for every $g\in\mathcal K_{\mathbf v}$ 
\begin{equation*}\begin{aligned} \sum_{k=1}^\infty   A_k^2\int_{\tau_k}|\varphi_1(\zeta)|^2|g(\zeta)|^2{\mathrm d}m(\zeta)  &
\leq \|g\|^2\sum_{k=1}^\infty A_k^2\int_{\tau_k}|\varphi_1(\zeta)|^2\|k_{\mathbf v,\zeta}\|^2{\mathrm d}m(\zeta) \\& 
\leq C_{n+1}^2\|g\|^2\sum_{k=1}^\infty A_k^2\int_{\tau_k}|\varphi_1(\zeta)|^2{\mathrm d}m(\zeta),\end{aligned}\end{equation*}
since $\tau_k\subset E_n$ for all $k\geq 1$. 

\emph{Second case}: $\delta_n>0$ for every $n\geq 1$.   Then there exists a subsequence $\{\delta_{n_k}\}_{k=1}^\infty$ such that 
 $0<\delta_{n_{k+1}}<\delta_{n_k}$ for all $k\geq 1$ and $\delta_{n_k}\to 0$. 
Set $$\tau_{1k}=\{\zeta\in E_{n_k}\ : \ |\varphi_0(\zeta)||\varphi_1(\zeta)|\leq\delta_{n_{k-1}}\}, \ \ \ k\geq 1.$$
Then $m(\tau_{1k})>0$ for every $k\geq 1$. Set $A_k=1/\delta_{n_{k-1}}$, $k\geq 1$. 
 Take a sequence $\{\varepsilon_k\}_{k=1}^\infty$ such that 
$\varepsilon_k>0$ for all $k\geq 1$, and \eqref{aaepsilon} is fulfilled. 
There  exists a sequence $\{\tau_k\}_{k=1}^\infty$ such that $\tau_k\subset\tau_{1k}$, $m(\tau_k)>0$, 
and $$C_{n_k+1}^2\int_{\tau_k}|\varphi_1|^2{\mathrm d}m\leq\varepsilon_k \ \text{ for all }k\geq 1.$$
Then \eqref{aavarphi1}  is fulfilled. 
Let $g\in\mathcal K_{\mathbf v}$. Since $\tau_k\subset E_{n_k}$, we have
$$ |g(\zeta)|=|(g,k_{\mathbf v,\zeta})|\leq\|g\|\|k_{\mathbf v,\zeta}\|\leq\|g\|C_{n_k+1}
\ \ \text{ for all }\zeta\in\tau_k.$$
Therefore, 
$$  \sum_{k=1}^\infty A_k^2\int_{\tau_k}|\varphi_1|^2|g|^2{\mathrm d}m
\leq \|g\|^2\sum_{k=1}^\infty C_{n_k+1}^2 A_k^2\int_{\tau_k}|\varphi_1|^2{\mathrm d}m.$$

In both cases, we obtain the sequences  $\{A_k\}_{k=1}^\infty$ and $\{\tau_k\}_{k=1}^\infty$ such that $A_k>0$ for all $k\geq 1$, 
$A_k\to\infty$, $m(\tau_k)>0$, $\tau_k\cap\tau_l=\emptyset$ for $k\neq l$, $k$, $l\geq 1$, \eqref{aavarphi1}  is fulfilled,
\begin{equation}\label{aakdeltak} \sup_{k\geq 1}A_k\mathop{\text{\rm ess\,sup}}_{\tau_k}|\varphi_0||\varphi_1|\leq 1, 
\end{equation}
and there exists $C>0$ such that 
\begin{equation}\label{aavarphi1g} \sum_{k=1}^\infty A_k^2\int_{\tau_k}|\varphi_1|^2|g|^2{\mathrm d}m\leq C\|g\|^2
\ \text{ for every }  g\in\mathcal K_{\mathbf v}.
\end{equation}

There exists an outer function $\varphi\in H^2$ such that $$|\varphi|=\begin{cases} |\varphi_1| & m\text{-a.e. on } 
\mathbb T\setminus\cup_{k=1}^\infty \tau_k, \\ 
A_k|\varphi_1|&  m\text{-a.e. on } \tau_k, \ \ k\geq 1.\end{cases}$$ 
It is easy to see that $\varphi$ satisfies to the conclusion of the lemma. 
Indeed, $\varphi\in H^2$ by \eqref{aavarphi1}. 
The inclusion $\varphi\varphi_0\in H^\infty$ follows from   the inclusion $\varphi_0\varphi_1\in H^\infty$ and \eqref{aakdeltak}. Since $A_k\to\infty$ and  $m(\tau_k)>0$ for all $k\geq 1$, we have 
$\varphi/\varphi_1\not\in H^\infty$. 
Let $g\in\mathcal K_{\mathbf v}$. Then
\begin{equation*}\begin{aligned}\int_{\mathbb T}|\varphi|^2|g|^2{\mathrm d}m &
 =\int_{\mathbb T\setminus\cup_{k=1}^\infty \tau_k}|\varphi|^2|g|^2{\mathrm d}m+
\sum_{k=1}^\infty \int_{\tau_k}|\varphi|^2|g|^2{\mathrm d}m\\&=
\int_{\mathbb T\setminus\cup_{k=1}^\infty \tau_k}|\varphi_1|^2|g|^2{\mathrm d}m+
\sum_{k=1}^\infty A_k^2 \int_{\tau_k}|\varphi_1|^2|g|^2{\mathrm d}m.\end{aligned}\end{equation*}
The first summand is finite by the assumption on $\varphi_1$, while the second summand is finite by \eqref{aavarphi1g}.
\end{proof}

\end{document}
